\subsection{Object Probability Distribution Calculation}

This section explains how the object probability distribution needed by the planner is created. This is only calculated when requested by the planner and only for the currently best particle. The information from the object maps, the room type map and the occupancy maps are combined for the estimation of the object probability distribution. This is done in three steps
\begin{enumerate}
 \item Improvement of the room type probability distributions using also the object maps and the occupancy grid mapped
 \item Improvement of the object probability distribution using also the room type probability distributions
 \item Estimation of the object probability distribution in areas not seen by the camera
\end{enumerate}

\subsubsection{Improvement of the room type probability distributions using the object maps and the occupancy grid mapped}

Using the correlation between objects in an area and the room type in this area the information object maps is used to get a better estimate of the room type probability distributions. First the probability of an object being near is calculated for every cell from the object maps with
\begin{equation}
P(O_{area, C_i}^{OT}|I_{1:t},\mathbf{x}_{1:t})=1-\prod_{C \in Area} (1-P(O_{D,C}^{OT}|I_{1:t},\mathbf{x}_{1:t}))
\end{equation}
The size of the area is chosen approximately the size of the visible area in a typical image used for the creation of the room type-object-correlation explained in section \ref{ssec:roomtypeobjcorr}. Therefore the values from the room-type-object-correlation-matrix can be used directly.

Assuming independence between the object probabilities given the room type, the complete room type probability is
\begin{multline}
P(R_{complete,Cj}=RT|I_{1:t},\mathbf{x}_{1:t}) = \\ \frac{1}{Z}P(R_{Cj}=RT | I_{1:t}, \mathbf{x}_{1:t})\prod_{k=1}^{|\mathbf{OT}|}\frac{P(R_{O_k,Cj}=RT|I_{1:t},\mathbf{x}_{1:t})}{P(R_{complete,Cj}=RT)}
\end{multline}
with 
\begin{equation}
Z=\sum_{RT \in \mathbf{RT}}\prod_{k=1}^{|\mathbf{OT}|}P(R_{Cj}=RT | I_{1:t}, \mathbf{x}_{1:t})\frac{P(R_{O_k,Cj}=RT|I_{1:t},\mathbf{x}_{1:t})}{P(R_{complete,Cj}=RT)}
\end{equation}

$P(R_{O_k,Cj}=RT|I_{1:t},\mathbf{x}_{1:t})$ is the probability of room type $RT$ in cell $C_j$ based on the probability of an object of type $O_k$ being near the cell. This probability is calculated by
\begin{multline}
P(R_{O_k,Cj}=RT|I_{1:t},\mathbf{x}_{1:t}) = \\ \sum_{v \in \{false,true\}}P(R_{O_k,Cj}=RT,O_{area, C_i}^{O_k}=v|I_{1:t},\mathbf{x}_{1:t}) \approx \\ \sum_{v \in \{false,true\}} P(O_{area, C_i}^{O_k}=v|I_{1:t},\mathbf{x}_{1:t})P(R_{O_k,Cj}=RT|O_{area, C_i}^{O_k}=v)
\end{multline}
$P(R_{O_k,Cj}|O_{area, C_i}^{O_k})$ can be calculated using the room-type-object-correlation matrix with 
\begin{multline}
P(R_{O_k,Cj}=RT|O_{area, C_i}^{O_k}=OT)=\\ \frac{P(R_{O_k,Cj}=RT)P(O_{area, C_i}^{O_k}=OT|R_{O_k,Cj}=RT)}{P(O_{area, C_i}^{O_k}=OT)} = \\ \frac{P(R_{O_k,Cj}=RT)P(O_{area, C_i}^{O_k}=OT|R_{O_k,Cj}=RT)}{\sum_{r \in \mathbf{RT}}P(R_{O_k,Cj}=r)P(O_{area, C_i}^{O_k}=OT|R_{O_k,Cj}=r)}
\end{multline}
with $P(O_{area, C_i}^{O_k}=OT|R_{O_k,Cj}=RT)$ being the matrix elements and $P(R_{O_k,Cj}=r) = \frac{1}{|\mathbf{RT}|}$.


\subsubsection{Improvement of the object probability distribution using the room type probability distributions}

First the object probability distribution $P(O_{R,C_i}^{OT})$ for all cells and object types based on the room type is calculated with
\begin{equation}
P(O_{R,C_i}^{OT})=\sum_{r\in \mathbf{RT}}P(R_{complete,Cj}=RT|I_{1:t},\mathbf{x}_{1:t})P_C(_{R,C_i}^{OT}|r)
\end{equation}

$P_C(_{R,C_i}^{OT}|r)$ is the probability of an object in a cell given the room type of the cell $r$. The room type-object-correlation matrix values are for areas of the usual size visible in an image, therefore the probabilities per cell have to be estimated. This is done by multiplying the matrix values with the estimated fraction of space visible in an image being occupied by an object. $P(O_{R,gauss,C_i}^{OT})$ is then a Gaussian filtered version of $P(O_{R,C_i}^{OT})$, which is done to model uncertainty of the positions of objects in a room. 


The object probability distribution using information from the object mapping and the room type is calculated with
\begin{equation}
P(O_{R,C_i}^{OT}|I_{1:t},\mathbf{x}_{1:t})=\frac{1}{Z}P(O_{D,C_i}^{OT}|I_{1:t},\mathbf{x}_{1:t})P(O_{R,gauss,C_i}^{OT})
\end{equation}
with the normalization term
\begin{multline}
Z=P(O_{D,C_i}^{OT}|I_{1:t},\mathbf{x}_{1:t})P(O_{R,gauss,C_i}^{OT})+\\(1-P(O_{D,C_i}^{OT}|I_{1:t},\mathbf{x}_{1:t}))(1-P(O_{R,gauss,C_i}^{OT}))
\end{multline} 

\subsubsection{Estimation of the object probability distribution in areas not seen by the camera}

Areas, where no information is available, are assumed to be outside the room and therefore the object probability is set to zero. Also the probabilities of areas known to be outside the room, as they are behind a doorway, are set to zero. This leaves areas only seen by the laser scanner, which has a greater field of view, to be considered. The room type of these areas is estimated based on the room type of neighboring cells, with the impact decreasing with distance, and the average room type in the room. Similar to above the object probabilities in the cells for all object types is calculated. 